CUDA将线程组织成block和32线程warp；寄存器、共享内存与global memory服务不同
复用范围\cite{cuda_programming_guide}。JHQ主码扫描将相邻候选映射到相邻线程，
精排则将一个候选的坐标分给多个lane。更多block能隐藏延迟，却重复LUT与缓冲；
寄存器和共享内存限制驻留。这些成本推动第5章的映射选择，具体带宽或cache归因
还需吞吐之外的计数器。

选择须与评分共同设计。Faiss采用warp选择\cite{johnson_gpu}；
bitonic、radix和merge原语在状态与工作量上有不同取舍。
AIR top-$k$与GridSelect研究依赖工作负载的全局选择\cite{gpu_topk}。
JHQ将这些已有思路用于路由和候选保留。

CUDA Graph回放减少launch设置，capture与工作区变化则属于执行生命周期。
共享界准入还需第5章的排序及完整key访问契约。标量分数正确或设备侧剖析
均不足以建立端到端正确性及主机可见性能。
